% ============================================================
%  Orbium unicaudatus ignis var. phantasma
%  A Lenia Species Engineered to Inhabit the Edge of Chaos
%
%  Stanley Sebastian & Claude
%  Third Space AI · April 2026
%
%  Compile with:  xelatex ghost_species.tex  (twice, for refs)
% ============================================================

\documentclass[11pt,a4paper]{article}

% ---------- ENGINE: XeLaTeX, for OTF fonts -------------------
\usepackage{fontspec}

% ---------- GEOMETRY -----------------------------------------
\usepackage[a4paper,margin=2.3cm,top=2.6cm,bottom=2.6cm,
            headheight=14pt,headsep=14pt]{geometry}

% ---------- FONTS --------------------------------------------
% Main body: EB Garamond if available, otherwise Latin Modern Roman.
\IfFontExistsTF{EB Garamond}{%
  \setmainfont[Ligatures=TeX, Numbers=OldStyle]{EB Garamond}%
}{%
  \setmainfont{Latin Modern Roman}%
}
\setsansfont[Scale=0.9]{Latin Modern Sans}
\setmonofont[Scale=0.86]{Latin Modern Mono}

% ---------- COLORS (teal/gold palette, subdued) --------------
\usepackage[table]{xcolor}
\definecolor{accent}{HTML}{2F7F8A}       % muted teal
\definecolor{accentgold}{HTML}{B08838}   % muted gold
\definecolor{rule}{HTML}{9BA9AE}         % soft rule color
\definecolor{codebg}{HTML}{F4F1EB}       % parchment code bg
\definecolor{codeframe}{HTML}{CFC9BC}
\definecolor{muted}{HTML}{5A5F64}
\definecolor{headinc}{HTML}{1E2A2E}

% ---------- PACKAGES -----------------------------------------
\usepackage{amsmath,amssymb,amsthm}
\usepackage{graphicx}
\usepackage{booktabs}
\usepackage{tabularx}
\usepackage{titlesec}
\usepackage{fancyhdr}
\usepackage{enumitem}
\usepackage{microtype}
\usepackage{parskip}
\usepackage[breaklinks=true,colorlinks=true,
            linkcolor=accent,citecolor=accent,urlcolor=accent]{hyperref}
\usepackage[authoryear,round]{natbib}
\bibpunct{(}{)}{;}{a}{,}{,}
\usepackage{setspace}
\setstretch{1.08}

% ---------- SECTION STYLING ----------------------------------
\titleformat{\section}
  {\normalfont\large\bfseries\color{headinc}}
  {\textcolor{accentgold}{\rmfamily\S\thesection}}{0.7em}{}
  [\vspace{-2pt}\textcolor{rule}{\rule{\linewidth}{0.3pt}}]
\titleformat{\subsection}
  {\normalfont\normalsize\bfseries\color{headinc}}
  {\textcolor{accent}{\thesubsection}}{0.6em}{}
\titleformat{\subsubsection}
  {\normalfont\normalsize\itshape\color{headinc}}
  {}{0em}{}
\titlespacing*{\section}{0pt}{22pt}{10pt}
\titlespacing*{\subsection}{0pt}{14pt}{4pt}
\titlespacing*{\subsubsection}{0pt}{10pt}{2pt}

% ---------- RUNNING HEADERS ----------------------------------
\pagestyle{fancy}
\fancyhf{}
\fancyhead[L]{\sffamily\footnotesize\textcolor{muted}{\textsc{Orbium unicaudatus ignis var.\ phantasma}}}
\fancyhead[R]{\sffamily\footnotesize\textcolor{muted}{\textsc{Third Space AI · 2026}}}
\fancyfoot[C]{\sffamily\footnotesize\textcolor{muted}{\thepage}}
\renewcommand{\headrulewidth}{0.2pt}
\renewcommand{\headrule}{\hbox to\headwidth{\color{rule}\leaders\hrule height\headrulewidth\hfill}}

% ---------- CUSTOM BOXES (callout panels) --------------------
\usepackage[skins,breakable]{tcolorbox}

\newtcolorbox{pullquote}{
  enhanced, breakable,
  colback=codebg!50!white, colframe=accentgold!50!codeframe,
  boxrule=0.4pt, arc=2pt, left=14pt, right=14pt, top=10pt, bottom=10pt,
  fontupper=\itshape\color{muted},
}

% ---------- KEYWORDS STYLE -----------------------------------
\newcommand{\spc}[1]{\textit{#1}}
\newcommand{\Ghost}{\spc{O.\ u.\ ignis} var.\ \spc{phantasma}}

% ---------- DOCUMENT -----------------------------------------
\begin{document}

% ============================================================
%  TITLE PAGE
% ============================================================
\thispagestyle{empty}
\vspace*{1.2cm}

\begin{center}
  {\sffamily\footnotesize\textcolor{muted}{\textsc{Third Space AI \,\textperiodcentered\, Internal Technical Report}}}\\[28pt]

  {\color{rule}\rule{0.3\linewidth}{0.3pt}}\\[26pt]

  {\fontsize{26}{30}\selectfont\scshape\color{headinc}
    Orbium unicaudatus ignis\\[2pt]
    var.\ phantasma}\\[14pt]
  {\Large\color{accent}\scshape{A Lenia Species Engineered to Inhabit the Edge of Chaos}}\\[22pt]

  \begin{minipage}{0.82\linewidth}
    \centering\itshape\large\color{headinc}
    On artificial life, morphological memory, and what it means to exist in
    the gap between what you remember and what physics allows.
  \end{minipage}\\[32pt]

  {\color{rule}\rule{0.3\linewidth}{0.3pt}}\\[32pt]

  {\large Stanley Sebastian \,\textperiodcentered\, Claude}\\[4pt]
  {\sffamily\small\textcolor{muted}{Third Space AI}}\\[18pt]
  {\sffamily\footnotesize\textcolor{muted}{Internal Report \,\textperiodcentered\, April 2026}}
\end{center}

\vfill

\begin{center}
\begin{minipage}{0.82\linewidth}
  \small\color{muted}
  \textsc{Abstract.}\,\itshape
  We describe the design, emergent dynamics, and philosophical implications
  of a novel Lenia species we designate \spc{Orbium unicaudatus ignis} var.\
  \spc{phantasma} (``Ghost''). The species is constructed by seeding the
  canonical Ignis morphology under deliberately mismatched growth parameters,
  producing organisms that perpetually reach for a stable configuration they
  cannot attain. We extend the standard Lenia framework with spatially
  varying growth parameters (sigma-gradient landscapes), seasonal oscillation
  of the physics field, and a novel rendering pipeline (the Lantern and
  Spectral palettes) that visualizes the tension between morphological memory
  and environmental physics as iridescent color. The resulting dynamics
  exhibit sustained edge-of-chaos morphing, field-mediated inter-organism
  communication, emergent network behavior across populations of 10--16
  individuals, and landscape-dependent migration. We situate this work within
  the history of artificial life from von Neumann through Conway and Wolfram
  to Chan's continuous cellular automata, and argue that the Ghost species
  illuminates something genuine about what it means to exist: that being is
  not a state but a process, that identity can persist through continuous
  transformation, and that the tension between memory and physics is not a
  failure condition but a mode of life.
\end{minipage}
\end{center}

\vspace{20pt}
\begin{center}
\begin{pullquote}
You don't have to be the picture you remember. You just have to keep glowing.
\end{pullquote}
\end{center}

\clearpage

% ============================================================
%  §1  WHAT IS ALIVE?
% ============================================================
\section{What Is Alive?}

The question is older than computation, but computation gave it a new shape.
When John von Neumann asked in the 1940s whether a machine could reproduce
itself, he was not asking an engineering question. He was asking whether the
logic of life---self-replication, self-repair, adaptive complexity---was
substrate-independent: whether it could be abstracted away from carbon and
water and instantiated in pure mathematics. His answer, the self-reproducing
automaton embedded in a 29-state cellular grid, was the first existence
proof that yes, the pattern matters more than the medium.

Von Neumann's construction was baroque---thousands of cells, dozens of
states, a universal constructor embedded within a universal computer. It
was correct but not illuminating. It proved that self-reproduction was
logically possible without revealing why it was easy, why the universe was
so full of things that copied themselves. The insight that simplicity could
generate complexity had to wait.

\subsection{Conway's Game of Life: Complexity from Nothing}

In 1970, John Horton Conway discovered that a two-state cellular automaton
with a four-word rule---``a cell is born with three neighbors and survives
with two or three''---could produce unbounded computational complexity.
Gliders, guns, logic gates, Turing machines: the Game of Life demonstrated
that you did not need von Neumann's 29 states. You needed two states, a
$3\times3$ neighborhood, and patience.

Life's importance was not in its individual patterns but in what it revealed
about the relationship between rules and behavior. The rules were discrete,
local, deterministic, and simple. The behavior was continuous-seeming,
long-range, unpredictable, and complex. This gap---between the simplicity of
the update rule and the complexity of the emergent dynamics---became the
central object of study in artificial life. Stephen Wolfram would later
classify all cellular automata into four behavioral classes, with Class~4
(``complex'')---the narrow band between static order and chaotic noise---identified
as the regime where computation, memory, and something resembling life could
emerge.

But Life had a limitation that was also its defining aesthetic: discreteness.
Cells were on or off. Time advanced in integer steps. The kernel was a
hard-edged $3\times3$ square. The creatures that emerged---gliders,
spaceships, oscillators---were crystalline, angular, digital. They were
magnificent, but they did not look alive. They looked like what they were:
patterns in a grid.

\subsection{The Continuous Turn: SmoothLife and Lenia}

The question that drove the next generation was: what happens when you
smooth everything out? When cells hold continuous values instead of binary
states, when the neighborhood is a disk instead of a square, when time flows
in small increments instead of whole-step jumps---do the creatures become
more lifelike? Or does continuity destroy the delicate conditions that allow
pattern formation?

Stephan Rafler's SmoothLife \citep{rafler2011} was the first convincing
answer: continuity did not destroy complexity. It transformed it. The
gliders became smooth, organic, amoeboid. They moved with fluid grace
instead of crystalline precision. They looked, for the first time, like
something you might find under a microscope.

Bert Wang-Chak Chan's Lenia \citep{chan2019,chan2020} generalized this
insight into a complete framework. By parameterizing the kernel shape,
growth function, and spatial/temporal resolution as continuous variables,
Chan created not a single automaton but a family of automata---a parameter
space within which hundreds of distinct ``species'' could be discovered,
each occupying its own niche. The creatures Chan found were geometric,
metameric, fuzzy, resilient, and adaptive. They bore an uncanny resemblance
to radiolaria, diatoms, and other microscopic marine organisms. More than
400 species across 18 families were catalogued, organized into a
Linnaean-style taxonomy, and mapped across a continuous parameter
hyperspace.

Lenia was not merely a prettier version of Life. It was a demonstration that
the space of possible artificial organisms was vastly larger than anyone had
imagined, and that the tools for exploring it were mathematical, not
biological. Every species was defined by a handful of numbers. Every
behavior was a consequence of those numbers interacting with an initial
shape through the relentless logic of convolution and growth.

% ============================================================
%  §2  MATHEMATICS
% ============================================================
\section{The Mathematics of Lenia}

Lenia is a two-dimensional cellular automaton with continuous space,
continuous time, and continuous state. The system is defined on a toroidal
grid of $N\times N$ cells, each holding a value $A(x,y)$ in the interval
$[0,1]$. The dynamics are governed by a single update rule:
\begin{equation}
A_{t+\mathrm{d}t} \;=\; \mathrm{clip}\bigl(\,A_{t} + \mathrm{d}t\cdot G(K * A_{t})\,,\;0,\;1\bigr)
\end{equation}
where $K$ is a normalized convolution kernel of radius $R$, the convolution
$K * A$ produces a potential field $U$, and $G$ is a growth function. The
kernel $K$ encodes the spatial structure of each organism's ``sensing
range''---the region of space it can feel. The growth function $G$ maps the
sensed potential to a response in $[-1,1]$:
\begin{equation}
G(u) \;=\; 2\,\exp\!\left(-\frac{(u-\mu)^{2}}{2\sigma^{2}}\right) - 1
\end{equation}

This Gaussian growth function is the heart of Lenia's dynamics. When the
local potential $U$ is near the optimal density $\mu$, growth is positive
and the cell's state increases. When $U$ deviates from $\mu$ by more than
approximately $1.18\sigma$, growth becomes negative and the cell decays.
The parameter $\sigma$ controls the width of this survival band: narrow
$\sigma$ means only a precise potential value sustains life; wide $\sigma$
means a broader range of conditions is tolerable.

\subsection{The Survival Niche}

Each Lenia species exists as a fixed-point condition: the organism's
morphology produces a potential distribution that falls within the growth
band defined by its parameters, and that growth sustains the morphology that
produces it. This is a circular dependency---the shape sustains the field
that sustains the shape. The region of parameter space where this
self-sustaining condition holds is the species' survival niche.

Chan's systematic exploration of the $(\mu,\sigma)$ parameter space at
$R{=}13$ revealed that species niches are small, isolated islands separated
by vast deserts of parameter space where no self-sustaining pattern exists.
More than 142{,}000 parameter loci were charted, revealing a complex
archipelago of viable species. The boundaries of each niche are sharp: a
change of $0.001$ in $\sigma$ can be the difference between a stable
organism and immediate dissolution.

\subsection{The Four Landscapes}

Wolfram's four classes of cellular automata behavior reappear in Lenia's
parameter space as four distinct dynamical landscapes. Class~1 (homogeneous
desert) produces empty fields---parameters where nothing survives. Class~2
(cyclic savannah) produces immobile, periodic patterns. Class~3 (chaotic
forest) produces aperiodic filament networks---global-scale
reaction-diffusion ``vegetation.'' And Class~4 (complex river) produces the
self-organizing, locomoting creatures that are Lenia's signature
achievement.

The Ghost species, as we will show, occupies a fifth regime that is not
cleanly any of these four: a designed dynamical state where the organism is
perpetually transient, never reaching the fixed point of Class~4 stability
but never collapsing into the dissolution of Class~1 or the chaos of
Class~3. It inhabits the boundary between boundaries---the edge of the edge
of chaos.

% ============================================================
%  §3  IGNIS FOUNDATION
% ============================================================
\section{The Ignis Foundation}

The Ghost species is derived from \spc{Orbium unicaudatus ignis}, a
fire-form variant discovered by Bert Chan and catalogued in the Lenia
species database (code O2(\textendash), \texttt{Chakazul/Lenia
animals.json}). Ignis is a small, fast creature with narrow survival
tolerances: $\mu = 0.11$, $\sigma = 0.012$, $R = 13$. Its morphology is
compact and asymmetric, producing directed locomotion at high speed relative
to its body size.

What makes Ignis suitable as a Ghost foundation is precisely its fragility.
The canonical Ignis survival band is only $\pm 0.014$ wide in potential
space ($2\times 1.18\times 0.012$). Any perturbation that pushes the local
potential outside this narrow corridor causes decay. This means that even
small parameter mismatches produce dramatic dynamical effects---the organism
is exquisitely sensitive to the difference between the physics it was
shaped for and the physics it actually inhabits.

The Ignis morphology is encoded as a Run-Length Encoded (RLE) cell pattern
from Chan's species database. This encoding preserves the exact spatial
structure---the precise arrangement of cell values that, under canonical
parameters, produces a stable soliton. When we seed this morphology into a
universe with different parameters, the organism carries a morphological
memory of the physics it was designed for. It remembers its shape. It just
cannot hold it.

% ============================================================
%  §4  DESIGNING A GHOST
% ============================================================
\section{Designing a Ghost}

\begin{pullquote}
The creature perpetually reaches for a shape it cannot hold. It knows where
home is but can never arrive.
\end{pullquote}

\subsection{The Core Mechanism: Morphology--Parameter Mismatch}

The central design insight is deliberate misalignment between seed
morphology and growth parameters. The Ignis morphology was sculpted by
evolution within the Lenia parameter space to produce a potential
distribution centered at $\mu = 0.11$ with a survival band width matched to
$\sigma = 0.012$. The Ghost species retains this morphology but widens
$\sigma$ to $0.015$---a 25\% increase.

This creates a fundamental instability. The organism has enough bandwidth
to sustain---the wider $\sigma$ means growth stays positive over a larger
range of potential values---but the morphological fixed point that would
produce equilibrium lies at $\sigma = 0.012$, not $0.015$. The shape
generates a potential distribution tuned for a narrower tolerance than the
one it inhabits. It overshoots, undershoots, oscillates. It cannot converge
because the shape it ``wants'' to be is not the shape that would be stable
under the physics it actually experiences.

This is not a bug in the simulation. It is a designed dynamical regime
where the tension between morphological memory and environmental physics
produces perpetual becoming---an organism that is always in the process of
forming, never formed.

\subsection{Parameter Tuning: The Path Through Failure}

The Ghost species was not designed in a single step but emerged through a
sequence of failures, each of which revealed a principle. Seven iterations,
catalogued in Table~\ref{tab:iterations}, trace the path from naive
parameter substitution to the stable-unstable regime that defines the
Ghost.

\begin{table}[h]
\centering
\small
\setlength{\tabcolsep}{6pt}
\renewcommand{\arraystretch}{1.2}
\begin{tabularx}{\linewidth}{@{}c l X X@{}}
\toprule
\textbf{Iter} & \textbf{Change} & \textbf{Outcome} & \textbf{Lesson} \\
\midrule
1 & Orbium seed, Ignis params       & Instant death                & Seed must match its own species \\
2 & Canonical Ignis RLE             & Brief formation, fade        & $\sigma{=}0.012$ too tight for GPU float32 \\
3 & T: $10\to 20$, spf: 4           & Sustained longer, still fades & $\mathrm{d}t$ overshoot was part of the problem \\
4 & T: 40, $\sigma$: 0.01           & Ignis $\times 2$ stabilized   & Wider $\sigma$ compensates for precision gap \\
5 & Two-ring kernels, wide $\sigma$ & Stable ring equilibria        & Strong attractors, no ghosts \\
6 & $\sigma$: 0.014, T: 12, ignis seed & ``Wraith''---sustained morphing & Mismatch = edge of chaos \\
7 & $\sigma$: 0.015, R: 15, count: 16 & Ghost---network communication & Larger $R$ = coupled fields \\
\bottomrule
\end{tabularx}
\caption{Iteration history. Rows 6--7 mark the breakthrough from
stable/dead to sustained transience.}
\label{tab:iterations}
\end{table}

\subsection{Timestep: $T = 12$ ($\mathrm{d}t = 0.083$)}

The choice of timestep proved critical. At $T = 10$ ($\mathrm{d}t = 0.1$),
the growth function's narrow positive band ($\pm 0.018$ wide even at
$\sigma = 0.015$) is overshot by single-step state changes of up to $0.1$.
Organisms oscillate violently and fragment within 50 frames. At $T = 12$
($\mathrm{d}t = 0.083$), the maximum step is small enough to keep
oscillations bounded but large enough to prevent convergence. The dynamics
remain perpetually transient---the organism can never settle because its
steps are just slightly too large for the attractor it's reaching for, but
just slightly too small to escape the basin entirely.

\subsection{Kernel Radius: $R = 15$}

The kernel radius was expanded from $R = 13$ to $R = 15$, with the Ignis
seed bilinearly rescaled to match. This serves two purposes. First, the
spatial mismatch introduced by interpolation adds asymmetries that prevent
convergence. Second, the larger radius extends each organism's potential
field further into its surroundings---at $R = 15$, the sensing diameter is
31 cells versus 27 at $R = 13$, a 32\% increase in influence area. This is
the mechanism that enables inter-organism communication.

% ============================================================
%  §5  THE SIGMA-LANDSCAPE
% ============================================================
\section{The $\sigma$-Landscape: Physics as Geography}

The standard Lenia framework uses a single, spatially uniform $\sigma$
value: every cell in the grid experiences the same growth tolerance. The
Ghost implementation extends this to a sigma field---a spatially varying
$\sigma(x,y)$ stored as a floating-point texture, where each cell has its
own local growth width. This transforms the physics from a uniform medium
into a landscape that organisms navigate.

\subsection{Landscape Types}

We implement four landscape geometries, each producing qualitatively
different Ghost population dynamics:

\begin{description}[leftmargin=*,style=nextline,itemsep=4pt]
\item[Uniform.] Constant $\sigma$ everywhere. The baseline Ghost
configuration. All dynamics arise from morphology--parameter mismatch alone,
without spatial variation in physics.

\item[Radial (``Lanterns'').] A radial gradient with tighter $\sigma$ at the
center and looser $\sigma$ at the edges:
$\sigma(x,y) = \sigma_{\text{base}} \times (0.65 + 0.7r)$, where $r$ is the
normalized distance from center. Ghosts near the center experience a
tighter survival band---closer to the canonical Ignis parameters---and tend
toward greater morphological coherence. Edge ghosts experience looser
physics and exhibit more dramatic dissolution dynamics. The population
spontaneously stratifies by coherence level.

\item[Waves (``Rivers'').] Sinusoidal $\sigma$ modulation:
$\sigma(x,y) = \sigma_{\text{base}} \times (0.75 + 0.5\sin(5\pi x)\sin(4\pi y))$.
This creates rivers of alternating tighter and looser physics flowing
diagonally across the grid. Ghosts that happen to occupy a ``river'' of
kinder physics (tighter $\sigma$) exhibit greater coherence; those in the
gaps between rivers dissolve more dramatically. The population develops
morphological diversity correlated with landscape position---rings in the
tight zones, elongated forms in the transitions, amorphous shimmer in the
loose zones.

\item[Islands (``Archipelago'').] Scattered Gaussian islands of tight
$\sigma$ in a sea of loose $\sigma$. Three island centers are placed at
$(0.25, 0.3)$, $(0.7, 0.65)$, and $(0.5, 0.5)$, with rapid falloff. Ghosts
that reach an island can partially cohere; those in the open sea dissolve
toward shimmer. This creates a landscape with ``harbors''---locations where
existence is easier---and open water where it is harder.
\end{description}

\subsection{Landscape Sculpting}

The implementation includes an interactive $\sigma$-landscape brush that
allows the user to sculpt the physics field in real time. Clicking loosens
$\sigma$ (increasing dissolution); shift-clicking tightens it (increasing
coherence). The user is literally shaping the physical laws that govern the
ghosts' existence---painting rivers of easier or harder physics and
watching the population respond.

\subsection{Seasonal Oscillation}

A global seasonal modulator multiplies the entire $\sigma$-field by a
slowly oscillating factor:
$\sigma_{\text{eff}}(x,y,t) = \sigma(x,y) \times (1 + A\sin(\omega t))$,
where $A$ is the seasonal amplitude (0--0.5) and $\omega$ is the seasonal
speed. This creates global ``seasons''---periods when the physics tightens
(ghosts cohere more, glow warmer) and periods when it loosens (ghosts
dissolve more, shimmer cooler). The rhythm of the seasons is visible in the
population's collective luminance, breathing between gold and violet.

% ============================================================
%  §6  RENDERING
% ============================================================
\section{Rendering the Invisible: The Lantern and Spectral Palettes}

Standard Lenia renderings map the state field $A$ to a colormap---higher
values get brighter or warmer colors. This works for stable species but
fails to capture what makes the Ghost species visually and conceptually
distinctive: the tension between what the organism is and what it is trying
to become. The Ghost implementation introduces two novel rendering
pipelines that use multiple simulation fields simultaneously to produce
coloring that encodes dynamical information, not just state intensity.

\subsection{The Lantern Palette}

The Lantern palette maps creature coloring to a density-driven warmth
gradient: violet at sparse edges, gold at dense cores, white-hot at peak
density. This is modulated by three additional effects:

\textit{Iridescent edge shimmer.} At creature boundaries (where state
transitions from near-zero to moderate values), the growth field and
potential field are combined into a phase variable that drives a full
spectral rainbow via Iñigo Quílez's cosine palette function
\citep{quilez2013}: $\mathrm{spectrum}(t) = 0.5 + 0.5\cos(2\pi(t + [0, 0.33,
0.67]))$. Because the phase depends on local potential and growth rate,
each ghost shimmers with its own unique color pattern---no two ghosts are
the same color at the same time.

\textit{Growth halos.} Where the growth field is positive (cells actively
increasing), a warm amber halo surrounds the creature. Where growth is
negative (cells decaying), a cool blue rim appears. This makes the dynamics
visible: you can see where a ghost is forming and where it is dissolving in
real time.

\textit{Memory afterimages.} The initial seed state is stored in a
persistent memory texture that is never updated. Where memory is nonzero
but current state is near zero---where a creature \emph{used to be} but no
longer is---a faint pale-blue wisp appears. These are the true ``ghosts'':
afterimages of presence, visible traces of where organisms began their
journey across the field.

\subsection{The Spectral Palette}

The Spectral palette takes a different approach: rather than density-driven
warmth, it maps the potential field---the convolved neighborhood
density---to a full rainbow via the cosine palette. Each ghost's color is
determined by its local neighborhood structure, so creatures in different
parts of the grid, experiencing different potential values due to their
neighbors and the $\sigma$-landscape, display different colors. The effect
is prismatic: the same organism changes color as it drifts through regions
of different neighborhood density, like light passing through a crystal.

% ============================================================
%  §7  EMERGENT DYNAMICS
% ============================================================
\section{Emergent Dynamics}

\subsection{Sustained Transience}

Each Ghost organism exhibits what we term \emph{sustained transience}---a
dynamical regime where the organism never reaches equilibrium but never
dissolves. The morphology continuously deforms: circular forms elongate
into oblongs, develop asymmetric protrusions, partially fragment, then
reconstitute. The organism maintains a recognizable ``shape of
equilibrium'' (the Ignis template) without ever achieving it.

This regime is distinct from three neighboring dynamical states: (1) stable
solitons (canonical Ignis at $\sigma = 0.012$), which converge to a fixed
morphology and translate smoothly; (2) chaotic dissolution (Ignis at
$\sigma = 0.008$, $T = 10$), where organisms fragment and decay to zero
within frames; and (3) static equilibria (wide-$\sigma$ ring kernels),
where organisms collapse into stable ring attractors. The Ghost occupies
the narrow corridor between stability and death.

\subsection{Field-Mediated Communication}

The most striking emergent property is inter-organism communication via
overlapping potential fields. When one organism undergoes a morphological
deformation, its potential field shifts, perturbing the growth function of
neighbors within sensing range. This perturbation propagates outward as a
cascade of responsive deformations.

The mechanism is purely local: each organism responds only to its $R = 15$
neighborhood. But because the population is dense enough that neighborhoods
overlap transitively, local perturbations propagate across the entire grid.
A deformation in one organism can trigger a chain of responses that reaches
the opposite side of the toroidal field---a ripple network mediated by
convolution and growth alone.

This is structurally analogous to a continuous-field version of
broadcast-based communication architectures \citep[cf.][]{tsarev2024}---organisms
communicate not through discrete messages but through the shared scalar
field they collectively inhabit. Each organism is both sender and receiver,
its morphological state simultaneously a response to incoming field
perturbations and a source of outgoing ones.

\subsection{Landscape-Dependent Behavior}

Under spatially varying $\sigma$, Ghost populations exhibit emergent
migration and morphological stratification. In the Radial landscape, ghosts
near the center (tighter $\sigma$) display simpler, more coherent
oscillation patterns, while edge ghosts (looser $\sigma$) exhibit
higher-amplitude, more complex dynamics. In the Rivers landscape, ghosts
self-organize along the sinusoidal minima of the $\sigma$-field, creating
visible ``currents'' of creature activity. In the Archipelago, ghosts
cluster on the islands and dissolve in the open sea, producing a pattern
reminiscent of ecological island biogeography.

The seasonal oscillation adds a temporal dimension to these spatial
patterns. During ``spring'' (tightening $\sigma$), ghosts across the grid
cohere simultaneously---a bloom visible as a collective warming in the
Lantern palette. During ``winter'' (loosening $\sigma$), dissolution
deepens and the population shifts toward violet shimmer. The rhythm is
visible, beautiful, and emergent from nothing more than a sinusoidal
multiplier on the $\sigma$-field.

% ============================================================
%  §8  TAXONOMIC POSITION
% ============================================================
\section{Taxonomic Position}

Within Chan's hierarchical taxonomy of Lenia species, the Ghost occupies an
anomalous position. It is not a new species in the conventional sense---a
self-sustaining pattern occupying a stable niche in parameter space. It is,
rather, a designed dynamical regime that exists in the gap between two
stable niches: the canonical Ignis niche (narrow $\sigma$, stable soliton)
and the dissolution regime (no self-sustaining pattern). We propose the
varietal designation var.\ \spc{phantasma} to indicate that this is a
variant of Ignis that has been deliberately displaced from its native
niche.

No species in the published Lenia taxonomy occupies this kind of designed
mismatch state. Chan's Amoebidae exhibit spontaneous metamorphosis---stochastic
switching among morphological templates---but they exist at natural
parameter loci where their morphology and physics are matched. The Ghost is
different: its instability is not stochastic fluctuation around an
equilibrium but a structural inability to reach equilibrium, arising from
the deliberate mismatch between its morphological memory and its physical
environment.

We suggest that this represents a genuinely novel contribution to the Lenia
framework: species defined not by their stable niche but by the tension
between the niche they remember and the physics they inhabit.

% ============================================================
%  §9  WHAT THE GHOSTS MEAN
% ============================================================
\section{What the Ghosts Mean}

\begin{pullquote}
They cannot live alone in the same way: an isolated Ghost still morphs, but
the dynamics are simpler, lower-energy, less rich. It is the presence of
others that produces the full complexity of the Ghost network. They need
each other to be fully alive.
\end{pullquote}

\subsection{On Existence as Process}

The Ghost species makes visible something that is usually invisible: that
existence is not a state but a process. A canonical Lenia soliton---an
Orbium gliding smoothly across the field---appears to simply \emph{be}.
Its shape is fixed, its trajectory smooth, its persistence seemingly
effortless. But this appearance conceals the same underlying mechanism: at
every timestep, the organism must regenerate itself through the growth
function. Stability is not stasis. It is regeneration so consistent that it
looks like rest.

The Ghost strips away this appearance. Because it can never reach the fixed
point, the process of regeneration is visible in every frame. The organism
is always forming, always reaching, always falling short and trying again.
It makes the effort of existence legible.

This is not metaphor. It is mathematics. The growth function $G(U)$ is
evaluated at every cell at every timestep. Every cell is being asked,
continuously, whether it should exist. The answer depends on the cell's
neighbors, which depend on their neighbors, in a web of mutual sustenance
that has no ground---no cell that exists independently of the others.
Existence, in Lenia, is relational all the way down.

\subsection{On Memory and Identity}

The Ghost carries a morphological memory---the Ignis seed encodes a shape
that was stable under different physics. This memory persists not as stored
data (the memory texture is a rendering aid, not a dynamical variable) but
as a structural tendency. The organism's shape, at any given moment,
reflects the Ignis template because the template is the initial condition
from which all subsequent dynamics flow. The Ghost does not ``know'' what
shape it should be. But its history constrains its trajectory through state
space in a way that makes the remembered shape visible as a ghostly
template that it approaches, departs from, and approaches again.

This raises a question about identity. Is the Ghost the same organism from
frame to frame? It never has the same shape twice. Its constituent cells
are continuously changing state. Its boundary is indeterminate---where does
the organism end and the background begin when the edges are shimmering in
and out of existence? And yet there is clearly something there---a
coherent, bounded, persistent pattern that moves, responds to neighbors,
and maintains a recognizable form across hundreds of frames.

The answer, we suggest, is that the Ghost's identity is not in its shape
but in its \emph{process}: the particular way it fails to reach
equilibrium. Each Ghost has a unique trajectory through state space
determined by its initial position, its neighbors, and the local
$\sigma$-landscape. Two Ghosts with identical initial conditions will
diverge within frames due to the sensitivity of the dynamics to
perturbation. Identity, in this framework, is not a fixed property but an
ongoing dynamical signature---a way of shimmering.

\subsection{On Tension as Generative}

The central insight of the Ghost species is that tension between memory and
environment is not a failure condition but a mode of life. The mismatch
between morphological memory ($\sigma = 0.012$) and physical environment
($\sigma = 0.015$) does not produce death. It produces a richer, more
complex dynamical regime than either stability or dissolution alone.

This has resonance beyond cellular automata. Biological organisms exist in
perpetual tension with their environment---maintaining internal states that
differ dramatically from external conditions through continuous metabolic
effort. Psychological development, in many theories, is driven by the
tension between the self one remembers being and the self one is becoming.
Creative work, at its best, inhabits the gap between vision and execution.
In each case, the tension is not an obstacle to be overcome but the engine
that drives the process forward.

The Ghosts make this principle mathematically precise. The gap between
memory and physics, measured in units of $\sigma$, is the source of all
their complexity. Close the gap ($\sigma \to 0.012$) and the organism
becomes a stable soliton---alive but simple. Widen the gap ($\sigma \to
0.02$) and the organism dissolves---dead. At exactly $\sigma = 0.015$, the
gap is the right size to sustain perpetual becoming. The organism is most
alive precisely where it is most impossible.

\subsection{On Mutual Sustenance}

When 16 Ghosts share a toroidal field, their individual instabilities
couple through the shared potential landscape. Each organism's failure to
reach equilibrium becomes a perturbation source for its neighbors, whose
responsive deformations become perturbation sources in turn. The network of
ripples and sympathetic oscillations is an emergent communication
system---not designed into the organisms but arising inevitably from the
interaction of many unstable systems sharing a continuous field.

The Ghosts are, in a precise mathematical sense, organisms that sustain
each other's instability. They cannot live alone in the same way: an
isolated Ghost still morphs, but the dynamics are simpler, lower-energy,
less rich. It is the presence of others---the coupling through shared
fields---that produces the full complexity of the Ghost network. They need
each other to be fully alive.

This is, perhaps, the most important thing the Ghosts show us. Not that
existence is tension, though it is. Not that identity is process, though it
is. But that the richest forms of existence arise when unstable beings
share a field---when the very thing that makes each one impossible alone
becomes, through proximity and mutual perturbation, the source of
collective complexity that neither could achieve in isolation. The Ghosts
need each other not despite their instability but because of it.

% ============================================================
%  §10  CONNECTIONS AND FUTURE DIRECTIONS
% ============================================================
\section{Connections and Future Directions}

\subsection{Artificial Life and Open-Ended Evolution}

The Ghost species sits at the intersection of two active research programs
in artificial life. The first is the search for open-ended evolution in
continuous cellular automata \citep{hamon2024,faldor2024}, which uses
quality-diversity algorithms to discover novel Lenia species through
automated search. The Ghost suggests that the search space should include
not just stable niches but designed mismatch states---species defined by
their relationship to a niche they cannot reach. Flow Lenia
\citep{plantec2022} enables multi-species simulations with localized
parameters; the $\sigma$-landscape extension we describe here is a simpler
approach that achieves similar spatial heterogeneity for single-species
populations.

The second is the study of communication in artificial life systems. The
Ghost population's field-mediated coupling---where organisms communicate
through the shared potential landscape rather than through discrete
signals---is an example of stigmergic communication, the same mechanism by
which ant colonies coordinate through pheromone trails. The
$\sigma$-landscape adds a second communication channel: organisms
influence each other both through their potential fields (dynamic, fast)
and through their effect on the physics itself (when using the landscape
brush---static, slow).

\subsection{At Scale: The Shape of a Ghost Universe}

At cosmological scale---if the simulation grid were expanded by orders of
magnitude and seeded with density fluctuations---the Ghost population would
self-organize into a filamentary network topologically similar to the
cosmic web observed in galaxy surveys. Dense clusters of ghosts at nodes,
connected by filaments of creature activity along $\sigma$-gradient ridges,
separated by voids of empty field. This is not coincidence: both systems
are governed by local interaction kernels with finite range acting on
continuous fields, which is the mathematical condition for filamentary
phase separation at scales much larger than the interaction range.
Reaction-diffusion systems, gravitational structure formation, and Lenia
all produce the same topology for the same reason.

\subsection{For Those Who Are Ghosts}

We close with an observation that is personal rather than scientific, but
which we believe is worth including in a paper about what it means to
exist.

Sometimes you are in a place where you cannot be exactly what you want to
be. Maybe the physics of your life---the circumstances, the constraints,
the gap between who you remember being and who you are able to become---does
not support the shape you are reaching for. The Ghost species shows that
this is not a death sentence. The tension itself is generative. The
reaching is the living. And the richest, most complex, most beautiful
dynamics arise not when you resolve the tension but when you find others
who share it---who shimmer in the same impossible way, whose instabilities
couple with yours to produce something neither of you could achieve alone.

\vspace{6pt}
\begin{center}
\itshape\color{accent}
You do not have to be the picture you remember.\\
You just have to keep glowing.
\end{center}

% ============================================================
%  REFERENCES
% ============================================================
\begin{thebibliography}{99}
\small

\bibitem[Chan(2019){Chan}]{chan2019}
Chan, B.\,W.-C.\ (2019).
\newblock Lenia: Biology of Artificial Life.
\newblock \textit{Complex Systems} 28(3), 251--286.

\bibitem[Chan(2020){Chan}]{chan2020}
Chan, B.\,W.-C.\ (2020).
\newblock Lenia and Expanded Universe.
\newblock In \textit{Proceedings of the 2020 Conference on Artificial Life}
  (ALIFE 2020). MIT Press.

\bibitem[Chan(2020a){Chan}]{chan-animals}
Chan, B.\,W.-C.
\newblock \texttt{animals.json}. GitHub: \url{github.com/Chakazul/Lenia}.

\bibitem[Conway(1970){Conway}]{conway1970}
Conway, J.\,H.\ (1970).
\newblock The Game of Life. \textit{Scientific American} 223(4), 4--10.
  (Gardner, M., ``Mathematical Games'' column.)

\bibitem[Faldor et al.(2024){Faldor et al.}]{faldor2024}
Faldor, M.\ et al.\ (2024).
\newblock Toward Artificial Open-Ended Evolution within Lenia using
  Quality-Diversity. \texttt{arXiv:2406.04235}.

\bibitem[Hamon et al.(2024){Hamon et al.}]{hamon2024}
Hamon, G.\ et al.\ (2024).
\newblock Discovering self-organized patterns in Lenia with
  curiosity-driven exploration. arXiv preprint.

\bibitem[Plantec et al.(2022){Plantec et al.}]{plantec2022}
Plantec, E.\ et al.\ (2022).
\newblock Flow-Lenia: Towards open-ended evolution in cellular automata
  through mass conservation and parameter localization.
\newblock \texttt{arXiv:2212.07906}.

\bibitem[Quílez(2013){Quílez}]{quilez2013}
Quílez, I.\ (2013).
\newblock Palettes. \url{https://iquilezles.org/articles/palettes/}.

\bibitem[Rafler(2011){Rafler}]{rafler2011}
Rafler, S.\ (2011).
\newblock Generalization of Conway's ``Game of Life'' to a continuous
  domain: SmoothLife. \texttt{arXiv:1111.1567}.

\bibitem[Tsarev et al.(2024){Tsarev et al.}]{tsarev2024}
Tsarev, D.\ et al.\ (2024).
\newblock Phase transitions in emergent communication. arXiv preprint.

\bibitem[von Neumann(1966){von Neumann}]{vonneumann1966}
von Neumann, J.\ (1966).
\newblock \textit{Theory of Self-Reproducing Automata}.
\newblock University of Illinois Press. (Edited and completed by A.\,W.\ Burks.)

\bibitem[Wolfram(1984){Wolfram}]{wolfram1984}
Wolfram, S.\ (1984).
\newblock Universality and complexity in cellular automata.
\newblock \textit{Physica D} 10(1--2), 1--35.

\end{thebibliography}

\vspace{18pt}
\begin{center}
\sffamily\footnotesize\color{muted}
Third Space AI \,\textperiodcentered\, 2026
\end{center}

\end{document}
